\fichatitulo{39 · RISCOS}{Frontiers in AI · 2025}{Alucinações: prompts e modelos}
{\small Anh-Hoang; Tran; Nguyen. \textit{Survey and Analysis of Hallucinations in Large Language Models: Attribution to Prompting Strategies or Model Behavior}. Frontiers in AI · 2025. \href{https://www.frontiersin.org/journals/artificial-intelligence/articles/10.3389/frai.2025.1622292/pdf}{Fonte consultada}.\par}

\campo{Problema e objetivo}
Como distinguir variação associada ao prompt e ao modelo nas alucinações?

\campo{Principal contribuição · síntese interpretativa}
Propõe um quadro diagnóstico com sensibilidade a prompts e variabilidade entre modelos, articulado a uma revisão de alucinações.

\campo{Método / natureza da evidência}
Descreve comparações entre modelos e variantes de instrução, com protocolos e métricas de atribuição.

\campo{Resultado ou argumento principal}
O enquadramento sugere examinar conjuntamente fatores de modelo e prompt, em vez de atribuir todo erro a um único componente.

\campo{Excertos-chave · seleção literal}
\begin{itemize}
\excerto{1}{Prompt Sensitivity (PS)}{Seção 4.2.}
\excerto{2}{Model Variability (MV)}{Seção 4.2.}
\end{itemize}

\campo{Limitações e análise crítica}
Análise da ficha: variação observada não basta para identificar causa. Há menções a GPT-4 na seção 4 e declaração de uso exclusivo de modelos abertos na seção 5; a relação entre exemplos e experimentos requer esclarecimento antes de reutilizar números.

\campo{Aproveitamento na dissertação · proposta}
Metodologia: testar sensibilidade a prompts; não adotar resultados quantitativos sem conferir o protocolo completo.

\registro{Fonte consultada nesta preparação: Resumo, definições, seção 4 e trechos da seção 5; consulta parcial. Chave: \texttt{anhHoang2025hallucinations}.}
